%% protocol_instances.tex
%% Paper F -- PCSS Theory: Protocol Instantiations (§4)
%% Issue #82 (F4)
%%
%% Instantiates DC(P), SBC analysis, and attack prediction for:
%%   Instance 1: LoRaMesher
%%   Instance 2: Meshtastic
%%   Instance 3: RPL
%%
%% Depends on: pcss_definitions.tex (Definitions 1--6), theorem1_proof.tex
%% Referenced in: SP-PaperF.tex §4

% ---------------------------------------------------------------------------
\section{Protocol Instantiations}
\label{sec:instances}
% ---------------------------------------------------------------------------

This section instantiates the PCSS framework (§3) on three deployed IoT mesh
routing protocols.
For each protocol we identify the Decision-Critical Set $\mathrm{DC}(P)$
(Definition~4), assess whether the Semantic Binding Condition $\mathrm{SBC}(P)$
holds in the default deployment (Definition~6), and derive the predicted attack
from the violated SBC condition.
Table~\ref{tab:instances-summary} summarises the results; §\ref{sec:taxonomy-corollary}
states the attack taxonomy corollary.

\begin{table}[t]
\centering
\caption{PCSS instantiation summary for three IoT mesh routing protocols.
  A \ding{55} in the SBC column means at least one condition of Definition~6
  fails in the default (unpatched) deployment.
  ``Patch target'' is the minimal message set that must receive
  an unforgeable authenticator to restore SBC.}
\label{tab:instances-summary}
\begin{tabular}{llcll}
\toprule
Protocol        & DC$(P)$                                   & SBC? & Attack category          & Patch target \\
\midrule
LoRaMesher      & \{Hello\}                                 & \ding{55} & Silent black hole       & HMAC on Hello \\
Meshtastic v2.6 & \{\texttt{MeshPacket.next\_hop}\}         & \ding{55} & Active relay (intercept) & Bind \texttt{next\_hop} in AEAD AAD \\
RPL (RFC 6550)  & \{DIO rank, DODAG version\}               & \ding{55} & DODAG subversion / reset & Authenticate DIO \\
\bottomrule
\end{tabular}
\end{table}

% ---------------------------------------------------------------------------
\subsection{Instance 1: LoRaMesher}
\label{sec:instances-lm}
% ---------------------------------------------------------------------------

\subsubsection*{Routing Decision Function}

In LoRaMesher~\cite{loramesher2023}, the distance-vector routing function for
node $i$ computes:
\[
  D_{r}^{i} = f(K_{i}, C_{ij}, Q_{i}, T_{i})
  \;:\;
  \mathrm{dst} \;\mapsto\; (\mathrm{next\_hop},\, \mathrm{metric})
\]
where the metric is the hop-count advertised in Hello messages, and the
next-hop selection rule is:
\[
  \mathrm{next\_hop}(\mathrm{dst}) \;=\;
  \arg\min_{n \in \mathrm{Neighbours}} \bigl(\mathrm{metric}_{n}(\mathrm{dst}) + 1\bigr).
\]

\subsubsection*{Decision-Critical Set}

\begin{proposition}
$\mathrm{DC}(P_{\mathrm{LM}}) = \{\texttt{Hello}\}$.
\end{proposition}

\begin{proof}
A Hello message carries the sender's node-id, a destination list, and
per-destination metric values.
Receiving a Hello is sufficient to update every routing entry for every
destination advertised; it is also necessary, since $D_{r}^{i}$ is updated only
on Hello receipt (not on data packets).
Hence Hello is both in $\mathrm{Cl}_{d}(\mathrm{Sem}(P_{\mathrm{LM}}))$ and
minimal: no strict subset is sufficient.
\end{proof}

\subsubsection*{SBC Violation}

In the baseline LoRaMesher implementation, Hello messages carry \emph{no
authenticator}: any node (or attacker) can broadcast a Hello with arbitrary
metric values.
Specifically, Definition~6 part~(1) fails: Hello $\in \mathrm{DC}(P_{\mathrm{LM}})$
but $\sigma_{m}$ is absent.
No condition of Definition~6 prevents an unauthorized party from producing a
valid-looking Hello.
Therefore $\mathrm{SBC}(P_{\mathrm{LM}})$ does not hold in the default deployment.

\subsubsection*{Attack Prediction and Field-Collapse Consequence}

By the SBC violation, an attacker $A$ can inject a Hello with metric~$= 0$
for a target destination dst.
This causes $D_{r}^{i}$ to install $A$ as the next-hop for dst:
\[
  \mathrm{RT}_{i}[\mathrm{dst}] \;\gets\; (\underbrace{A}_{\mathrm{next\_hop}},\; 0).
\]

However, LoRaMesher employs a \emph{dst~$=$~next\_hop field structure}: the
packet header sets \texttt{dst} equal to \texttt{next\_hop} at the MAC layer.
When a data packet arrives at $A$, the MAC header reads
\texttt{dst} $=$ \texttt{next\_hop} $= A$, so the packet is delivered
\emph{to $A$ as the final destination} and is not forwarded.
Since $A$ is not the intended application-layer destination, the packet is
silently discarded: \textbf{silent black hole}.

The field-collapse is formally captured by the domain structure of $D_{r}^{i}$:
the routing table maps dst $\to$ next\_hop, but the MAC layer enforces
dst $\equiv$ next\_hop, collapsing the image of $D_{r}^{i}$ to singletons.
This is distinct from protocols where dst and next\_hop are independent fields.

\subsubsection*{Hardware Evidence}

Paper~B experiment E5b: under the spoofing attack, PDR drops to $0\%$ across
all tested topologies (linear, mesh, star).
After applying HMAC-SHA256 on Hello (the minimal patch that restores
$\mathrm{SBC}(P_{\mathrm{LM}})$), PDR is restored to baseline.
The Tamarin lemma \texttt{L1\_MetricAuthenticity} is falsified on the
baseline model and verified on the patched model (\texttt{patched\_lora\_dv\_2node.spthy}).

\subsubsection*{Patch and SBC Restoration}

Adding HMAC-SHA256 to every Hello message with key $K$ shared by
$\mathrm{AuthorizedSet}(P_{\mathrm{LM}})$ restores $\mathrm{SBC}(P_{\mathrm{LM}})$:
all three conditions of Definition~6 are satisfied.
Theorem~\ref{thm:sbc-sufficiency} then applies, and the silent black-hole
attack is impossible (Corollary~\ref{cor:bc-impossible}).
The physical cost of HMAC-SHA256 per Hello is ${\approx}2\,\mathrm{ms}$ CPU and
$32$ bytes airtime overhead; at one Hello per 10 s, this costs
${\approx}0.17\,\mathrm{mAh/day}$ on EoRa~PI (ESP32-S3, 1000~mAh), well within
the feasibility envelope (Theorem~2, §5).

% ---------------------------------------------------------------------------
\subsection{Instance 2: Meshtastic}
\label{sec:instances-mesh}
% ---------------------------------------------------------------------------

\subsubsection*{Routing Decision Function}

In Meshtastic v2.6+~\cite{meshtastic2024}, the routing function selects the
next relay for a packet via the \texttt{MeshPacket.next\_hop} field in the
packet header.
NodeInfo messages carry node identity and are authenticated via AEAD
(AES-256-CTR with a channel key), but \emph{they do not determine the
forwarding path}: forwarding is determined solely by \texttt{next\_hop}.

\subsubsection*{Decision-Critical Set}

\begin{proposition}
$\mathrm{DC}(P_{\mathrm{Mesh}}) = \{\texttt{MeshPacket.next\_hop}\}$.
\end{proposition}

\begin{proof}
The relay selection at node $i$ is:
\[
  D_{r}^{i} \;:\; \mathrm{dst} \;\mapsto\; \mathrm{MeshPacket.next\_hop}.
\]
Modifying \texttt{next\_hop} directly changes the forwarding decision without
altering any other header field; NodeInfo messages influence the
neighbor table but not the per-packet forwarding decision.
Hence $\mathrm{DC}(P_{\mathrm{Mesh}}) = \{\texttt{MeshPacket.next\_hop}\}$,
and NodeInfo $\notin \mathrm{DC}(P_{\mathrm{Mesh}})$.
\end{proof}

\subsubsection*{SBC Violation}

Although NodeInfo is AEAD-authenticated, \texttt{MeshPacket.next\_hop} is
\emph{advisory} in the Meshtastic packet format: it is carried in the
\textbf{unencrypted} portion of the header and is \emph{not} included in the
AEAD AAD.
Definition~6 part~(1) therefore fails for $\mathrm{DC}(P_{\mathrm{Mesh}})$:
the authenticator on NodeInfo does not bind \texttt{next\_hop}.
An attacker who can observe or relay a packet can overwrite \texttt{next\_hop}
with its own node-id without breaking any cryptographic tag.

\begin{remark}[HMAC on NodeInfo Does Not Fix This]
A common but incorrect patch intuition is to add authentication to NodeInfo
messages.
This does not restore $\mathrm{SBC}(P_{\mathrm{Mesh}})$ because NodeInfo
$\notin \mathrm{DC}(P_{\mathrm{Mesh}})$: it does not determine the forwarding
decision.
The correct patch is to include \texttt{next\_hop} in the AEAD AAD of
\texttt{MeshPacket}, binding the relay identity to the per-packet authenticator.
\end{remark}

\subsubsection*{Attack Prediction and Active-Relay Consequence}

An attacker $A$ rewrites \texttt{MeshPacket.next\_hop} to its own node-id in a
packet destined for node $D$.
The MAC layer at the preceding hop forwards the packet to $A$.
Unlike LoRaMesher, the Meshtastic header carries \textbf{separate} dst and
next\_hop fields: \texttt{dst} remains $D$, so $A$ recognises that the packet
must be forwarded and delivers it (optionally after copying the payload).
The consequence is an \textbf{active relay attack}: $A$ becomes an
intermediate relay, achieving 100\% payload intercept.

The structural difference from Instance~1 is captured by the domain of $D_{r}^{i}$:
in $P_{\mathrm{Mesh}}$, $D_{r}^{i}$ maps dst $\to$ next\_hop with
dst $\neq$ next\_hop in general.
The image of $D_{r}^{i}$ is not collapsed to singletons, so the attacker
receives and can forward the packet, rather than silently dropping it.

\subsubsection*{Hardware Evidence}

Paper~C experiment E16: PDR $= 100\%$ and intercept $= 100\%$ under the
unpatched protocol across all tested topologies.
Both metrics drop to $\leq 5\%$ (noise floor) when \texttt{next\_hop} is
bound in AEAD AAD.
The Tamarin lemma \texttt{L3\_No\_Relay\_Takeover} is falsified on the
baseline model and verified on the patched model
(\texttt{meshtastic\_patch.spthy}).

% ---------------------------------------------------------------------------
\subsection{Instance 3: RPL}
\label{sec:instances-rpl}
% ---------------------------------------------------------------------------

\subsubsection*{Routing Decision Function}

In RPL (RFC~6550)~\cite{rfc6550}, the DODAG routing function for node $i$ is:
\[
  D_{r}^{i} \;:\; \mathrm{dst} \;\mapsto\; \mathrm{preferred\_parent}
\]
where preferred\_parent is selected by minimising rank among received DIO
(DODAG Information Object) messages.
The DODAG version number is used to trigger global topology resets.

\subsubsection*{Decision-Critical Set}

\begin{proposition}
$\mathrm{DC}(P_{\mathrm{RPL}}) = \{\texttt{DIO.rank},\; \texttt{DIO.version}\}$.
\end{proposition}

\begin{proof}
Rank determines preferred\_parent selection: a DIO advertising rank $r$
causes any node with current rank $> r + \mathrm{MinHopRankIncrease}$ to
adopt the sender as preferred parent, updating $D_{r}^{i}$.
DODAG version number triggers a DODAG reset: any node receiving a DIO with
a higher version number MUST rejoin and discard its current routing table.
Both fields are therefore sufficient and necessary to control the routing
topology; DAO (destination advertisement) messages are downstream of
preferred\_parent selection and do not independently control it.
\end{proof}

\subsubsection*{SBC Violation}

RFC~6550 does not mandate DIO authentication in the default profile; the
optional RPL security mode (Section~10 of RFC~6550) specifies a security
object but it is disabled by default in major implementations
(Contiki-NG, RIOT OS).
Definition~6 part~(1) fails: DIO messages carrying rank and version are
unauthenticated, so any on-path or in-range node can forge or replay them.

\subsubsection*{Attack Predictions}

Two independent SBC violations lead to two distinct attack categories:

\begin{itemize}
  \item \textbf{Rank manipulation.}\;
    An attacker broadcasts a DIO with rank $= 1$ (minimal),
    causing all neighbouring nodes to adopt it as preferred parent.
    The DODAG is subverted into a topology where all traffic flows through $A$.
    Literature: Raoof et al.~\cite{raoof2019rpl} catalogue this as the
    ``rank attack''; Perrey et al.~\cite{perrey2016trail} propose the TRAIL
    mechanism as a countermeasure, confirming the prediction.

  \item \textbf{Version attack.}\;
    An attacker broadcasts a DIO with a future version number $v' > v_{\mathrm{current}}$,
    forcing a global DODAG reset.
    All nodes must rejoin, discarding routing state and triggering a
    convergence storm.
    This is a denial-of-service attack derived directly from the SBC
    violation on \texttt{DIO.version}.
    Literature: \cite{raoof2019rpl,dvir2011vera} confirm that unauthenticated
    version increments cause topology oscillation.
\end{itemize}

Unlike Instances~1 and~2, RPL has no hardware evidence in this paper series;
the prediction is validated by the independent literature above.

% ---------------------------------------------------------------------------
\subsection{Attack Taxonomy Corollary}
\label{sec:taxonomy-corollary}
% ---------------------------------------------------------------------------

\begin{corollary}[DC$(P)$ Structure Determines Attack Category]
\label{cor:taxonomy}
Let $P_{1}$ and $P_{2}$ be two protocols with the same type of SBC violation
(Definition~6 part~(1) fails), but different structures of $D_{r}^{i}$:
\begin{itemize}
  \item[$P_{1}$:] $D_{r}^{i}$ has dst $\equiv$ next\_hop
    (field collapse; $\mathrm{Im}(D_{r}^{i})$ maps each dst to itself).
  \item[$P_{2}$:] $D_{r}^{i}$ has dst $\neq$ next\_hop in general
    (separate fields; $\mathrm{Im}(D_{r}^{i})$ is unconstrained).
\end{itemize}
Then an attacker exploiting the SBC violation in $P_{1}$ produces a
\emph{silent black-hole} attack, whereas the same SBC violation in $P_{2}$
produces an \emph{active-relay} attack.
\end{corollary}

\begin{proof}
In $P_{1}$, forwarding a packet to next\_hop $= A$ delivers the packet to
$A$ as the final MAC-layer destination.
Since dst $\equiv$ next\_hop, the packet header gives $A$ no information about
the application-layer destination $D$; $A$ discards the packet (silent black hole).
In $P_{2}$, forwarding to next\_hop $= A$ delivers the packet to $A$ with
dst $= D \neq A$ visible in the header.
$A$ recognises it is an intermediate relay, copies the payload, and
(optionally) forwards, achieving active relay with full intercept.
The SBC violation --- inability to forge a valid authenticator for
$m \in \mathrm{DC}(P)$ without key $K$ --- is identical in both cases;
the attack category is determined entirely by the structural property of $D_{r}^{i}$.
\hfill$\square$
\end{proof}

\begin{remark}[Cross-Paper Recovery]
Corollary~\ref{cor:taxonomy} formally recovers Paper~C's header-semantics
taxonomy~\cite{paperC} as a corollary of the PCSS framework.
Paper~C introduced the dst/next\_hop field-collapse observation empirically;
the PCSS framework provides the structural explanation: it is a consequence of
$\mathrm{Im}(D_{r}^{i})$ collapsing to singletons in $P_{\mathrm{LM}}$ but not
in $P_{\mathrm{Mesh}}$.
Table~\ref{tab:taxonomy} summarises the cross-paper corollary.
\end{remark}

\begin{table}[t]
\centering
\caption{Attack taxonomy as corollary of PCSS.
  All three protocols share the same SBC violation type; attack category is
  determined by $D_{r}^{i}$ structure.}
\label{tab:taxonomy}
\begin{tabular}{llll}
\toprule
Protocol        & SBC violation         & $D_{r}^{i}$ structure & Attack category \\
\midrule
LoRaMesher      & Hello unauthenticated  & dst $\equiv$ next\_hop & Silent black hole (Paper~B) \\
Meshtastic      & next\_hop advisory     & dst $\neq$ next\_hop   & Active relay / intercept (Paper~C) \\
RPL             & DIO unauthenticated    & rank $\to$ parent      & DODAG subversion / reset \\
\bottomrule
\end{tabular}
\end{table}
